\chapter{Il profilo del giocatore di futsal}
\label{cap:profilogiocatore}

\textbf{Profilo del giocatore}: quadro delle qualità fisiche del calciatore a 5 ricavato dai
\textbf{test di valutazione} pubblicati in letteratura internazionale, usato come \textbf{punto di
riferimento} per classificare il singolo atleta.

L'argomento è diviso in tre parti. Si muove dai \textbf{dati descrittivi} --- antropometria e
massimo consumo di ossigeno --- e si prosegue con \textbf{studi longitudinali} che seguono una
squadra per un'intera stagione, così da avere non solo i valori ma anche la loro
\textbf{evoluzione nel tempo}; seguono l'\textbf{epidemiologia degli infortuni} e, in chiusura, il
profilo della \textbf{giocatrice}. Il dato di riferimento serve a rispondere a una domanda
operativa: sottoposto un giocatore a un \textit{counter movement jump} o a uno \textit{Yo-Yo
test}, quel risultato lo colloca o no fra i soggetti d'élite?

I valori qui riportati sono \textbf{medie}: la particolarità di ogni squadra e di ogni giocatore fa
sì che il singolo se ne discosti anche abbondantemente. Ciò che conta è il \textbf{trend}, non il
valore puntuale.

% ---------------------------------------------------------------------------
\section{I dati antropometrici}

Fonte: rassegna di \textbf{Matzenbacher e coll.} (2014), che raccoglie i dati di una trentina di
studi su giocatori professionisti, ordinati per autore e livello di gioco.

\rubrica{Le grandezze considerate}
\begin{itemize}
  \item \textbf{età};
  \item \textbf{peso} corporeo;
  \item \textbf{statura};
  \item \textbf{composizione corporea}, espressa come percentuale di massa grassa.
\end{itemize}

\rubrica{Il trend}
\begin{itemize}
  \item \textbf{peso}: mediamente fra \textbf{70 e 80\,kg};
  \item \textbf{statura}: il giocatore di futsal \textbf{non ha la struttura del calciatore},
    attestandosi su valori inferiori;
  \item \textbf{massa grassa}: fra l'\textbf{8 e il 12\,\%} della massa totale --- valore
    \textbf{in linea} con quello del calciatore.
\end{itemize}

\rubrica{Il caso della prima divisione spagnola}
Il dato della \textbf{1\textsuperscript{a} divisione spagnola} --- livello internazionale, studio
di \textbf{Jiménez e coll.} --- si discosta nettamente da tutti gli altri della rassegna. La
ragione è il contesto: in Spagna il campionato è \textbf{professionistico}, esattamente come la
Serie A, B e C del calcio in Italia. La tendenza che ne emerge è quella di giocatori
\textbf{sempre più strutturati} anche dal punto di vista antropometrico.

Lo stesso studio è l'unico della rassegna che separa i due ruoli:

\begin{itemize}
  \item \textbf{portieri} ($n = 3$): 27,6 $\pm$ 5 anni; 78,6 $\pm$ 6,5\,kg; 184 $\pm$ 2\,cm;
  \item \textbf{giocatori di movimento} ($n = 9$): 24,5 $\pm$ 3 anni; 76,5 $\pm$ 6,8\,kg;
    180 $\pm$ 12,3\,cm.
\end{itemize}

Il portiere è dunque \textbf{più anziano, più pesante e più alto} del giocatore di movimento.

% ---------------------------------------------------------------------------
\section{Il massimo consumo di ossigeno}

\textbf{$VO_2max$}: capacità di estrarre ossigeno dall'aria e di utilizzarlo nei processi
metabolici per ricavarne energia.

\rubrica{I valori della rassegna}
\begin{itemize}
  \item \textbf{alto livello}: valori tendenzialmente \textbf{sopra i 55\,ml/kg/min}, quindi
    \textbf{simili a quelli del calciatore};
  \item i valori più elevati vengono dalla \textbf{2\textsuperscript{a} divisione spagnola}:
    64,8 (53,8--75,8) e 65,1 $\pm$ 6,2\,ml/kg/min;
  \item il \textbf{livello regionale brasiliano} sta più in basso: 52,6 $\pm$ 3,1;
    55,7 $\pm$ 3,7; 59,6 $\pm$ 2,5\,ml/kg/min;
  \item la \textbf{3\textsuperscript{a} divisione italiana} si colloca a 55,2 $\pm$ 5,7\,ml/kg/min.
\end{itemize}

\rubrica{Portieri e giocatori di movimento}
Nella 1\textsuperscript{a} divisione brasiliana lo scarto fra i ruoli è marcato:
\begin{itemize}
  \item \textbf{portieri} ($n = 22$): 50,6 $\pm$ 5,24\,ml/kg/min;
  \item \textbf{giocatori di movimento} ($n = 164$): 59 $\pm$ 5,8\,ml/kg/min.
\end{itemize}

Un solo studio riporta il dato \textbf{prima e dopo} la stagione: 71,5 $\pm$ 5,9\,ml/kg/min in
entrata contro 67,6 $\pm$ 3,5 in uscita.

% ---------------------------------------------------------------------------
\section{La distribuzione del carico in una stagione}

Studio di \textbf{Miloski e coll.} (2015), \textit{Seasonal training load distribution of
professional futsal players: effects on physical fitness, muscle damage and hormonal status}.

\rubrica{Scopo}
Descrivere la \textbf{distribuzione del carico di lavoro} durante una stagione nel futsal
brasiliano di alto livello.

\rubrica{Metodo di indagine}
Studio \textbf{comparativo longitudinale} su una squadra di Serie A brasiliana, seguita dal
precampionato al termine del campionato.

\subsection{Il controllo del carico}

Tre famiglie di rilevazioni, ciascuna con una propria periodicità.

\begin{enumerate}
  \item \textbf{Test di valutazione} (siglati \textbf{PT}, \textit{physical tests}):
    \begin{itemize}
      \item \textit{counter movement jump} (CMJ): forza esplosiva e uso dell'energia elastica;
      \item \textbf{sprint da 5\,m}: capacità di accelerazione;
      \item \textbf{sprint da 20\,m}: profilo di velocità;
      \item \textbf{T-test}: agilità, capacità fondamentale nel calcio a 5; per il protocollo
        vedi cap.~\ref{cap:metabolismoanaerobico}.
    \end{itemize}
  \item \textbf{Carico interno percepito}: \textbf{RPE} con \textbf{scala di Borg CR10}
    (valori da 0 a 10), somministrata \textbf{a ogni seduta};
  \item \textbf{Componente ematica e ormonale} (siglata \textbf{BS}, \textit{blood sample}),
    rilevata in momenti determinati della stagione:
    \begin{itemize}
      \item \textbf{creatinchinasi} (CK): indicatore di \textbf{danno muscolare};
      \item \textbf{testosterone};
      \item \textbf{cortisolo}.
    \end{itemize}
\end{enumerate}

\subsection{La periodizzazione}

La stagione è divisa in \textbf{sei periodi di allenamento} su 22 settimane. Le sigle dei tipi di
lavoro sono:

\begin{itemize}
  \item \textbf{STG}: allenamento della \textbf{forza};
  \item \textbf{AER/L TOL}: allenamento \textbf{aerobico} e di \textbf{tolleranza al lattato};
  \item \textbf{SP/PW}: \textbf{velocità e potenza};
  \item \textbf{TT SKL}: lavoro \textbf{tecnico-tattico};
  \item \textbf{REC}: \textbf{recupero}.
\end{itemize}

\begin{table}[htbp]
  \centering
  \small
  \renewcommand{\arraystretch}{1.3}
  \begin{tabularx}{\textwidth}{|c|c|X|c|c|}
    \hline
    \textbf{Periodo} & \textbf{Settimane} & \textbf{Obiettivi principali} &
      \textbf{Sedute TT SKL} & \textbf{Partite (am./uff.)} \\
    \hline
    1 & 1--4 & AER/L TOL, STG & 23 & 2 / 0 \\
    \hline
    2 & 5--6 & SP/PW, TT SKL & 16 & 2 / 0 \\
    \hline
    3 & 7--10 & TT SKL, REC & 22 & 0 / 7 \\
    \hline
    4 & 11--13 & TT SKL, SP/PW & 25 & 0 / 1 \\
    \hline
    5 & 14--17 & TT SKL, SP/PW & 26 & 0 / 6 \\
    \hline
    6 & 18--22 & TT SKL, REC & 29 & 0 / 7 \\
    \hline
  \end{tabularx}
\end{table}

\rubrica{Che cosa si legge nella tabella}
\begin{itemize}
  \item la \textbf{preparazione precampionato} dura \textbf{sei settimane} e cura, con pesi
    diversi, \textbf{tutte le capacità condizionali}; vi si disputano \textbf{4 amichevoli};
  \item all'\textbf{inizio del campionato} il carico fisico-atletico è \textbf{alleggerito in modo
    importante} e cresce la quota tecnico-tattica;
  \item nelle \textbf{settimane 11--13}, periodo di \textbf{sosta}, lo staff \textbf{rialza} il
    carico fisico-atletico --- è l'unico momento della stagione in cui torna il lavoro aerobico e
    di tolleranza al lattato --- con 25 sedute e una sola partita ufficiale;
  \item nelle \textbf{ultime settimane} il lavoro torna in prevalenza tecnico-tattico, con
    \textbf{13 partite ufficiali} complessive fra i periodi 5 e 6.
\end{itemize}

\rubrica{I dettagli dei mezzi allenanti}
\begin{itemize}
  \item \textbf{STG}: presente in \textbf{ogni} periodo. Nel precampionato 4 serie, da 8--10 RM a
    4--6 RM, con recuperi fra le serie di 75--90\,s o di 180\,s nelle fasi a carico più alto; in
    stagione si stabilizza su 3 serie di 6--8 o 8--10 RM con 75--90\,s di recupero;
  \item \textbf{AER/L TOL}: 3 serie di 4$\times$60\,s alla \textbf{MAS} (\textit{maximal aerobic
    speed}) con 60\,s di cammino fra le prove, poi 3 serie di 3$\times$30\,s \textit{all out} con
    30\,s di cammino; 150\,s fra le serie. Compare \textbf{solo due volte} nell'anno: nelle
    \textbf{settimane 1--4} e di nuovo nella \textbf{sosta} (settimane 11--13, 3 sedute di
    2 serie di 5$\times$15\,s \textit{all out}, 45\,s di cammino, 120\,s fra le serie);
  \item \textbf{SP/PW}: da 1 serie di 8 sprint in linea retta si passa a
    4$\times$8 \textit{drop jumps}, 2$\times$6 sprint da 10\,m con \textit{traction belt} e
    2$\times$6 sprint da 5--20\,m; in stagione gli sprint diventano \textbf{con palla} e
    \textbf{con cambio di direzione}. Recuperi di 60--90\,s passivi;
  \item \textbf{TT SKL}: \textit{ball-drill}, \textit{small-sided games} e partite simulate, il cui
    scopo cambia secondo le necessità tecnico-tattiche della squadra.
\end{itemize}

\subsection{L'andamento settimanale del carico}

Il carico settimanale è espresso in \textbf{unità arbitrarie} e confrontato con quattro bande di
riferimento: \textit{low}, \textit{moderate-low}, \textit{moderate-high}, \textit{high}.

\begin{itemize}
  \item nelle \textbf{prime tre settimane} di preparazione il carico è ai valori
    \textbf{più alti} della stagione, sopra la banda \textit{high};
  \item segue un \textbf{brusco calo} alla quarta settimana, poi una risalita in banda
    \textit{moderate-high} nelle settimane 5--6;
  \item durante il \textbf{campionato} il carico oscilla fra le bande intermedie, con i picchi
    più alti in corrispondenza della \textbf{sosta};
  \item nel \textbf{finale di stagione} il trend è di \textbf{progressiva diminuzione}, fino al
    minimo assoluto dell'ultima settimana.
\end{itemize}

\subsection{L'evoluzione dei test fisici}

I quattro test sono stati somministrati \textbf{prima} della preparazione (PT$_1$), a
\textbf{fine} preparazione (PT$_2$), al momento della \textbf{sosta} (PT$_3$) e a
\textbf{fine campionato} (PT$_4$).

\begin{table}[htbp]
  \centering
  \small
  \renewcommand{\arraystretch}{1.3}
  \begin{tabularx}{\textwidth}{|X|c|c|c|c|}
    \hline
    \textbf{Test} & \textbf{PT$_1$} & \textbf{PT$_2$} & \textbf{PT$_3$} & \textbf{PT$_4$} \\
    \hline
    CMJ (cm) & 47,5 $\pm$ 5,5 & 47,8 $\pm$ 6,1 & 49,1 $\pm$ 6,2 & 49,8 $\pm$ 6,2 \\
    \hline
    Sprint 5\,m (s) & 1,10 $\pm$ 0,08 & 1,08 $\pm$ 0,05 & 1,04 $\pm$ 0,07 & 1,00 $\pm$ 0,04 \\
    \hline
    Sprint 20\,m (s) & 3,14 $\pm$ 0,11 & 3,09 $\pm$ 0,11 & 3,03 $\pm$ 0,13 & 3,00 $\pm$ 0,07 \\
    \hline
    T-test (s) & 9,24 $\pm$ 0,31 & 8,75 $\pm$ 0,30 & 8,71 $\pm$ 0,22 & 8,56 $\pm$ 0,22 \\
    \hline
    $VO_2max$ (ml/kg/min) & 49,5 $\pm$ 3,5 & 52,3 $\pm$ 3,7 & 53,4 $\pm$ 2,8 & 53,3 $\pm$ 2,9 \\
    \hline
  \end{tabularx}
\end{table}

\rubrica{Come vanno letti gli scarti}
Gli autori non usano il solo valore di $p$ ma la \textbf{\textit{magnitude-based inference}}, che
gradua la certezza della differenza rispetto a PT$_1$ su quattro livelli: \textbf{possibile}
(25--75\,\%), \textbf{probabile} (75--95\,\%), \textbf{molto probabile} (95--99\,\%),
\textbf{quasi certa} ($>$99\,\%).

\rubrica{Il trend}
\begin{itemize}
  \item \textbf{forza esplosiva} (CMJ): durante la preparazione \textbf{non aumenta} --- i
    giocatori non sono stati sollecitati con quell'obiettivo --- e cresce invece
    \textbf{durante il campionato}, con differenza probabile da PT$_3$ in poi;
  \item \textbf{sprint da 5 e da 20\,m}: stesso andamento, con miglioramento che diventa
    \textbf{quasi certo} a fine stagione. I due test si correlano con il CMJ, perché indagano la
    stessa capacità di forza esplosiva;
  \item \textbf{T-test}: migliora \textbf{già a fine preparazione} e continua a migliorare, con
    differenza \textbf{quasi certa} in tutti i rilievi successivi;
  \item \textbf{$VO_2max$}: sale \textbf{subito}, di quasi il 6\,\% già a fine preparazione, e poi
    si stabilizza.
\end{itemize}

Sullo sprint da 20\,m va ricordato che il calciatore a 5 \textbf{non percorre quasi mai} sprint
così lunghi: dalla match analysis (cap.~\ref{cap:matchanalysis}) lo sprint medio dello studio di
Castagna e D'Ottavio è di circa \textbf{10\,m}, e comunque si aggira fra i 10 e i 15\,m --- effetto
combinato della presenza di palla e avversario e della \textbf{grandezza del terreno di gioco}.

\subsection{Danno muscolare e profilo ormonale}

I prelievi (BS$_1$--BS$_7$) sono stati eseguiti \textbf{ogni due settimane} durante la
preparazione e in punti determinati del campionato.

\begin{table}[htbp]
  \centering
  \small
  \renewcommand{\arraystretch}{1.3}
  \begin{tabularx}{\textwidth}{|X|c|c|c|c|c|c|c|}
    \hline
    & \textbf{BS$_1$} & \textbf{BS$_2$} & \textbf{BS$_3$} & \textbf{BS$_4$} & \textbf{BS$_5$} &
      \textbf{BS$_6$} & \textbf{BS$_7$} \\
    \hline
    CK (U/l) & 156,8 & 266,3 & 246,9 & 179,3 & 200,4 & 201,1 & 215,6 \\
    \hline
    Testosterone (pg/ml) & 21,6 & 23,7 & 23,1 & 23,1 & 22,5 & 24,0 & 21,9 \\
    \hline
    Cortisolo ($\mu$g/dl) & 14,2 & 12,3 & 14,5 & 12,9 & 16,8 & 14,9 & 15,8 \\
    \hline
    Rapporto T:C & 1,5 & 2,0 & 1,7 & 1,9 & 1,4 & 1,7 & 1,4 \\
    \hline
  \end{tabularx}
\end{table}

\rubrica{La creatinchinasi}
\begin{itemize}
  \item valori \textbf{alti} seguono gli allenamenti fortemente sollecitanti;
  \item a \textbf{fine preparazione} il valore è nettamente più alto che all'inizio, e la
    differenza è \textbf{statisticamente significativa}: testimonia l'\textbf{ampio carico di
    lavoro} somministrato;
  \item nel corso della stagione il valore \textbf{scende e si assesta} intorno alle
    \textbf{200\,U/l}.
\end{itemize}

È il parametro su cui concentrarsi: dà un \textbf{feedback preciso} su volume e intensità del
lavoro effettivamente prodotto in campo, e quindi un'indicazione fondamentale per la
programmazione.

\rubrica{Il rapporto testosterone/cortisolo}
Resta \textbf{costante} per tutto l'anno. Per \textit{stress} non si intendono qui i soli aspetti
psicologici del periodo agonistico: la costanza del rapporto indica che a questo livello i
giocatori \textbf{assorbono} lo stress della stagione.

\subsection{Che cosa se ne ricava}

\begin{enumerate}
  \item a \textbf{fine preparazione} migliorano le qualità effettivamente sollecitate --- in
    particolare il \textbf{$VO_2max$}, con un guadagno del \textbf{5,7\,\%} --- segno di
    adattamento a un carico elevato con obiettivo specifico;
  \item le qualità \textbf{prettamente neuromuscolari} (CMJ, sprint da 5 e da 20\,m)
    \textbf{non} migliorano in modo significativo durante la preparazione: il miglioramento
    \textbf{avviene poi durante la stagione};
  \item la \textbf{CK}, dopo l'incremento del precampionato, ha un andamento \textbf{costante}:
    i giocatori hanno la \textbf{capacità di sostenere} il carico di lavoro e di gara,
    adattandosi in modo positivo;
  \item il \textbf{rapporto T:C} costante indica una grande capacità di sostenere \textbf{lunghi
    periodi di stress agonistico} --- ed è l'andamento da attendersi nell'alto livello;
  \item sulla \textbf{programmazione} --- endurance e forza nel precampionato, velocità e potenza
    durante la stagione --- lo studio \textbf{non dà una risposta generalizzabile}: descrive le
    scelte di un singolo staff, e altre scuole di pensiero collocherebbero diversamente le
    sollecitazioni. Di questo articolo vanno presi la \textbf{precisione} e la
    \textbf{dovizia di dati}, non il modello di periodizzazione;
  \item ne esce rafforzata l'\textbf{importanza del controllo e della valutazione}
    dell'allenamento: in una rosa di una quindicina di elementi la somministrazione del carico è
    \textbf{di squadra}, ma la \textbf{risposta è individuale}. Solo un'analisi neuromuscolare e
    metabolica di questo dettaglio permette di \textbf{individualizzare} la proposta.
\end{enumerate}

% ---------------------------------------------------------------------------
\section{L'evoluzione delle capacità fisiche nella stagione}

Studio di \textbf{Oliveira e coll.} (2013), \textit{Seasonal changes in physical performance and
heart rate variability in high level futsal players}. Dove Miloski descrive \textit{come è stato
distribuito} il carico, qui si guarda a \textit{che cosa cambia} nelle capacità fisiche, con il
fuoco sulla \textbf{preparazione precampionato}.

\rubrica{Scopo}
Determinare le modificazioni delle capacità fisiche \textbf{prima} e \textbf{durante} la stagione
sportiva.

\rubrica{Campione}
11 giocatori di alto livello: 24,3 $\pm$ 2,9 anni; 176,3 $\pm$ 5,2\,cm; 76,1 $\pm$ 6,3\,kg ---
valori \textbf{confrontabili} con quelli della rassegna antropometrica.

\subsection{I due test}

\begin{enumerate}
  \item \textbf{Yo-Yo intermittent recovery test livello 1} (Yo-Yo IR1): indaga la capacità di
    \textbf{sostenere momenti di alta intensità} e di \textbf{recuperare in fretta} per poterli
    ripetere. Si svolge a navetta fra due cinesini posti a 20\,m, con la partenza scandita da un
    \textbf{segnale acustico} che fa da lepre: il giocatore deve trovarsi al cinesino in coincidenza
    del \textit{bip}. Fra una navetta di 20 $+$ 20\,m e la successiva si inserisce un
    \textbf{recupero attivo} camminando 5\,m di andata e 5 di ritorno, in attesa del segnale che fa
    ripartire. La frequenza del \textit{bip} cresce, e con essa la velocità di corsa, fino
    all'\textbf{esaurimento}: richiede quindi il \textbf{massimo impegno} dell'atleta. Il risultato
    è la \textbf{distanza percorsa};
  \item \textbf{Repeated Sprint Ability} (RSA): 6 sprint da 40\,m, percorsi come navetta
    20 $+$ 20\,m con \textbf{cambio di direzione a 180\textdegree}, alla \textbf{massima velocità}
    e non a ritmo imposto, separati da \textbf{20\,s di recupero passivo} --- cioè da fermo.
\end{enumerate}

I momenti di rilevazione sono tre: \textbf{M1} all'inizio della preparazione precampionato,
\textbf{M2} alla fine della preparazione (che qui dura \textbf{tre settimane}), \textbf{M3} a metà
stagione.

\subsection{La settimana tipo}

Gli autori la documentano su due settimane campione: la \textbf{terza}, cioè l'ultima di
preparazione, e la \textbf{quinta}, già a campionato iniziato.

\begin{itemize}
  \item nel \textbf{precampionato} il carico fisico-atletico è alto sia per volume sia per
    intensità: potenza muscolare in palestra, \textit{interval training} a media intensità,
    \textit{repeated sprint}, forza muscolare, attività aerobiche a bassa intensità, con il
    tecnico-tattico relegato al pomeriggio e un'amichevole il sabato;
  \item \textbf{in stagione} crescono il lavoro \textbf{tecnico-tattico} e le \textbf{attività
    rigenerative}, e compaiono gli spostamenti per le trasferte;
  \item le partite si giocano il \textbf{lunedì} e il \textbf{giovedì} sera --- mediamente
    \textbf{due a settimana} --- per impegni di campionato e di coppa.
\end{itemize}

\subsection{I risultati}

\begin{table}[htbp]
  \centering
  \small
  \renewcommand{\arraystretch}{1.3}
  \begin{tabularx}{\textwidth}{|X|c|c|c|}
    \hline
    & \textbf{M1} & \textbf{M2} & \textbf{M3} \\
    \hline
    Yo-Yo IR1 (m) & 1244 (298) & 1491 (396) & 1465 (270) \\
    \hline
    RSA$_{mean}$ (s) & 7,43 (0,2) & 7,24 (0,2) & 7,13 (0,3) \\
    \hline
    RSA$_{best}$ (s) & 6,93 (0,2) & 6,89 (0,2) & 6,69 (0,3) \\
    \hline
    RSA$_{decrement}$ (\%) & 6,7 (0,3) & 5,0 (0,9) & 6,6 (2,4) \\
    \hline
  \end{tabularx}
\end{table}

Valori come media (deviazione standard). Gli autori accompagnano ogni confronto con l'\textbf{
\textit{effect size}} e il suo intervallo di confidenza al 95\,\%, graduato in
\textbf{trascurabile} ($<$0,2), \textbf{piccolo} (0,2--0,6), \textbf{moderato} (0,6--1,2) e
\textbf{grande} ($>$1,2).

\rubrica{Che cosa migliora nella preparazione}
\begin{itemize}
  \item \textbf{Yo-Yo IR1}: da poco più di 1200\,m a quasi 1500, differenza significativa rispetto
    a M1 ed \textit{effect size} \textbf{moderato}; il valore poi \textbf{si mantiene} per tutta la
    stagione (confronto M2 \textit{vs} M3 trascurabile);
  \item \textbf{RSA$_{decrement}$}: scende da 6,7 a 5,0\,\%, con \textit{effect size}
    \textbf{grande}. È il parametro che più interessa in un test RSA, perché misura proprio la
    capacità di \textbf{ripetere} lo sprint recuperando in fretta;
  \item \textbf{RSA$_{mean}$}: migliora già in preparazione e continua a migliorare.
\end{itemize}

\rubrica{Che cosa non migliora}
\begin{itemize}
  \item \textbf{RSA$_{best}$} resta praticamente invariato fra M1 e M2 (\textit{effect size}
    \textbf{trascurabile}) e migliora solo a metà stagione, dove la differenza è significativa
    rispetto a \textbf{entrambi} i rilievi precedenti: l'allenamento precampionato \textbf{non
    aveva come obiettivo} il miglioramento della velocità;
  \item \textbf{RSA$_{decrement}$} fa il percorso inverso: guadagnato in preparazione, a M3
    \textbf{torna quasi al valore di partenza} (6,6\,\%). La lettura proposta è che il carico
    somministrato in stagione \textbf{non fosse specifico} per questa caratteristica.
\end{itemize}

\rubrica{I riferimenti da portarsi in campo}
\begin{itemize}
  \item nello \textbf{Yo-Yo IR1} un giocatore d'élite percorre \textbf{oltre 1200\,m} a inizio
    stagione e attorno ai \textbf{1500\,m} una volta preparato;
  \item nel \textbf{decremento} dell'RSA ci si attende un valore \textbf{intorno al 7\,\%} a inizio
    stagione, migliorabile se la caratteristica viene sollecitata;
  \item il giocatore di calcio a 5 si rivela \textbf{molto sensibile} all'aumento di carico e
    \textbf{si adatta rapidamente}: sia la capacità di ripetere sprint sia quella di sostenere
    l'alta intensità prolungata migliorano in modo netto in sole tre settimane.
\end{itemize}

% ---------------------------------------------------------------------------
\section{L'epidemiologia degli infortuni}

Studio di \textbf{Martinez-Riaza e coll.} (2017), \textit{Epidemiology of injuries in the Spanish
national futsal male team: a five-season retrospective study}.

\rubrica{Scopo}
Analizzare l'\textbf{assistenza medica} fornita ai giocatori prima e durante le partite del
campionato spagnolo.

\rubrica{Metodo}
Studio \textbf{retrospettivo} su \textbf{5 stagioni} (dalla 2010-11 alla 2014-15), condotto
somministrando \textbf{questionari} allo staff medico e ai giocatori. Il totale è di
\textbf{411 prestazioni mediche}.

\subsection{La distribuzione per ruolo}

\begin{itemize}
  \item \textbf{laterale}: 50,4\,\% (207 interventi);
  \item \textbf{ultimo}: 20\,\% (82);
  \item \textbf{pivot}: 16,8\,\% (69);
  \item \textbf{portiere}: 10,2\,\% (42);
  \item \textbf{ala-pivot}: 2,7\,\% (11).
\end{itemize}

La lettura è tattica: gli interventi si concentrano sui ruoli che per compito
\textbf{percorrono più campo} e sono \textbf{più sollecitati nell'alta intensità} --- laterali
anzitutto, poi ultimi, che organizzano il gioco, e pivot, il riferimento offensivo. Restano in
basso portiere e ala-pivot; quest'ultimo ha caratteristiche offensive ma \textbf{non un'occupazione
fissa dello spazio} come il pivot.

\subsection{I meccanismi}

\begin{itemize}
  \item \textbf{intrinseci} --- tutto ciò che avviene \textbf{senza trauma esterno}, come le
    lesioni muscolari: \textbf{67,2\,\%};
  \item \textbf{estrinseci} --- veri e propri \textbf{urti}: intervento dell'avversario o cadute
    particolari: \textbf{30,2\,\%};
  \item \textbf{origine non conosciuta}: 2,7\,\%.
\end{itemize}

Il peso dei due meccanismi \textbf{cambia con il ruolo}: nei portieri e negli ultimi gli
estrinseci sfiorano il 45\,\%, mentre nei laterali scendono al 26,1\,\% e nei pivot al 17,4\,\%,
dove gli intrinseci arrivano all'81,2\,\%.

\subsection{La natura degli infortuni}

\begin{itemize}
  \item \textbf{57,7\,\%} degli interventi riguarda il \textbf{tessuto muscolare};
  \item il \textbf{sovraccarico} è di gran lunga la voce più numerosa: oltre il \textbf{50\,\%}
    dei casi;
  \item gli infortuni al \textbf{ginocchio} sono il \textbf{10\,\%} del totale, quelli alla
    \textbf{caviglia} il \textbf{6,5\,\%}: numeri \textbf{non elevati}, contro l'attesa comune;
  \item il \textbf{47\,\%} degli infortuni --- quasi la metà --- \textbf{avviene in allenamento},
    non in partita.
\end{itemize}

% ---------------------------------------------------------------------------
\section{Il profilo della giocatrice di futsal}

Il quadro fin qui tracciato riguarda il futsal maschile. Sul femminile la letteratura è
\textbf{molto più scarsa} --- come già si era visto per la match analysis
(cap.~\ref{cap:matchanalysis}) --- e i tre studi che seguono ne coprono altrettanti versanti:
il confronto fra livelli di qualificazione, la fitness aerobica, l'epidemiologia degli infortuni.

\subsection{Elite e sub-elite a confronto}

Studio di \textbf{Ramos-Campo e coll.} (2016), \textit{Physical performance of elite and subelite
Spanish female futsal players}. Doppio scopo: \textbf{valutare} le caratteristiche fisiche e
tecniche delle giocatrici e \textbf{confrontare} elite (D1) e sub-elite (D2).

\rubrica{I sei parametri}
\begin{itemize}
  \item \textbf{CMJ}: forza esplosiva e riuso dell'energia elastica;
  \item \textbf{SJ} (\textit{squat jump}): espressione più \textbf{pura} di forza esplosiva, senza
    contributo elastico;
  \item \textbf{sprint da 30\,m}: velocità e accelerazione, fortemente correlate ai due salti;
  \item \textbf{30\,m agility}: lo stesso sprint da 30\,m, ma percorso a \textbf{zigzag fra
    paletti}, cioè su un tracciato a esse invece che in linea;
  \item \textbf{Repeated Sprint Ability};
  \item \textbf{massima velocità di tiro}: unico parametro \textbf{tecnico}, misurato con un
    \textbf{radar}.
\end{itemize}

\begin{table}[htbp]
  \centering
  \small
  \renewcommand{\arraystretch}{1.3}
  \begin{tabularx}{\textwidth}{|X|c|c|c|c|}
    \hline
    & \multicolumn{2}{c|}{\textbf{CMJ}} & \multicolumn{2}{c|}{\textbf{SJ}} \\
    \cline{2-5}
    & $P_{max}$ (W/kg) & $h$ (cm) & $P_{max}$ (W/kg) & $h$ (cm) \\
    \hline
    Elite & 43,1 $\pm$ 4,2 & 26,7 $\pm$ 0,3 & 42,2 $\pm$ 2,4 & 26,1 $\pm$ 0,4 \\
    \hline
    Sub-elite & 43,6 $\pm$ 4,7 & 24,3 $\pm$ 0,3 & 40,8 $\pm$ 3,6 & 24,2 $\pm$ 0,3 \\
    \hline
    $p$ & 0,743 & 0,167 & 0,366 & 0,207 \\
    \hline
  \end{tabularx}
\end{table}

\begin{table}[htbp]
  \centering
  \small
  \renewcommand{\arraystretch}{1.3}
  \begin{tabularx}{\textwidth}{|X|c|c|c|c|c|}
    \hline
    & \textbf{30\,m} (s) & \textbf{RSA$_{best}$} (s) & \textbf{RSA$_{total}$} (s) &
      \textbf{RSA$_{change}$} (\%) & \textbf{RSA$_{dec}$} (\%) \\
    \hline
    Elite & 4,9 $\pm$ 0,2 & 4,9 $\pm$ 0,2 & 41,1 $\pm$ 1,9 & 3,6 $\pm$ 3,9 & 4,9 $\pm$ 2,2 \\
    \hline
    Sub-elite & 5,0 $\pm$ 0,2 & 5,0 $\pm$ 0,2 & 41,8 $\pm$ 1,7 & 4,6 $\pm$ 5,1 & 4,7 $\pm$ 1,8 \\
    \hline
    $p$ & 0,456 & 0,282 & 0,317 & 0,601 & 0,821 \\
    \hline
  \end{tabularx}
\end{table}

\rubrica{Dove i due livelli non si distinguono}
Su \textbf{nessuno} dei parametri fisico-atletici puri emerge una differenza statisticamente
significativa: né nei due salti, né sui 30\,m in linea, né nell'RSA --- dove il valore assoluto
dell'elite è appena migliore, mentre il \textbf{decremento} è addirittura leggermente peggiore
(4,9\,\% contro 4,7\,\%).

\rubrica{Dove invece si distinguono}
\begin{itemize}
  \item \textbf{30\,m agility}: 5,5\,s l'elite contro 5,7\,s la sub-elite, differenza
    significativa;
  \item \textbf{velocità di tiro}: 84,5\,km/h contro 79,8\,km/h, differenza significativa.
\end{itemize}

\rubrica{La riflessione}
Delle due l'una: o i test \textbf{non erano specifici} per il calcio a 5 --- ipotesi poco
sostenibile, visto il tipo di prove --- oppure il \textbf{livello di qualificazione non dipende}
dalle caratteristiche fisiche. A discriminare sono un
\textbf{parametro tecnico} (il tiro) e una \textbf{capacità specifica} del futsal (l'agility), non
la prestazione fisico-atletica pura. Come riferimento operativo: per collocarsi fra le giocatrici
d'élite occorre tirare \textbf{sopra gli 80\,km/h}.

\subsection{La fitness aerobica e il FIET}

Studio di \textbf{Barbero-Alvarez e coll.} (2015), \textit{Aerobic fitness and performance in
elite female futsal players}. Due test messi a confronto per \textbf{validare} una prova da campo
specifica del futsal contro il riferimento di laboratorio.

\rubrica{FIET (Futsal Intermittent Endurance Test)}
Step di corsa a navetta \textbf{3 $\times$ 15\,m} (45\,m), con la velocità dettata da un
\textbf{segnale acustico} a regime progressivo fino all'esaurimento; velocità iniziale
\textbf{9\,km/h}. Si distingue dallo Yo-Yo IR1, che è invece una navetta di 2 $\times$ 20\,m.
\begin{itemize}
  \item distanza totale: \textbf{1125,0 $\pm$ 121,0\,m};
  \item velocità di picco: 15,2 $\pm$ 0,5\,km/h;
  \item $FC_{picco}$: 199 $\pm$ 8\,bpm;
  \item lattato di picco: 12,5 $\pm$ 2,2\,mmol/l.
\end{itemize}

\rubrica{Test incrementale su treadmill}
Partenza a \textbf{6\,km/h}, con incrementi di \textbf{2\,km/h ogni 3 minuti} fino
all'esaurimento.
\begin{itemize}
  \item $VO_2max$: \textbf{45,3 $\pm$ 5,6\,ml/kg/min};
  \item velocità aerobica massima (MAS): 12,5 $\pm$ 1,77\,km/h;
  \item $FC_{max}$: 197 $\pm$ 8\,bpm;
  \item lattato di picco: 11,3 $\pm$ 1,4\,mmol/l.
\end{itemize}

\rubrica{Che cosa se ne ricava}
\begin{itemize}
  \item le giocatrici di futsal hanno livelli di fitness aerobica \textbf{moderati} rispetto alle
    calciatrici di pari livello ed età --- e nettamente inferiori ai 55\,ml/kg/min e oltre visti
    per il futsal maschile;
  \item i valori di lattato, sopra le 11\,mmol/l in \textbf{entrambe} le prove, mostrano un forte
    coinvolgimento del metabolismo \textbf{anaerobico lattacido}, e le frequenze cardiache vicine
    ai 200\,bpm confermano che l'impegno è stato massimale;
  \item il \textbf{FIET} può quindi considerarsi un \textbf{valido test da campo} per valutare la
    fitness aerobica e anaerobica nel futsal femminile.
\end{itemize}

\subsection{Gli infortuni nel futsal femminile}

Studio di \textbf{Gayardo e coll.} (2012), \textit{Prevalence of injuries in female athletes of
brazilian futsal}. Indagine con \textbf{questionario} su un campione di \textbf{147 atlete} fra i
\textbf{16 e i 35 anni}, per un totale di \textbf{104 infortuni} riportati.

\begin{table}[htbp]
  \centering
  \small
  \renewcommand{\arraystretch}{1.3}
  \begin{tabularx}{\textwidth}{|X|c|c|}
    \hline
    \textbf{Sede} & \textbf{N} & \textbf{\%} \\
    \hline
    Caviglia & 30 & 28,9 \\
    \hline
    Coscia & 25 & 24,0 \\
    \hline
    Ginocchio & 24 & 23,1 \\
    \hline
    Mano e polso & 8 & 7,7 \\
    \hline
    Piede & 7 & 6,7 \\
    \hline
    Gomito & 2 & 1,9 \\
    \hline
    Anca & 2 & 1,9 \\
    \hline
    Gamba & 2 & 1,9 \\
    \hline
    Altro & 4 & 3,9 \\
    \hline
    \textbf{Totale} & \textbf{104} & \textbf{100} \\
    \hline
  \end{tabularx}
\end{table}

\rubrica{La gravità}
Il criterio è il \textbf{tempo di degenza}: il periodo in cui lo staff medico ferma l'atleta, che
non si allena o si allena in forma \textbf{differenziata}, e comunque non gioca.
\begin{itemize}
  \item \textbf{lieve} (fino a 6 giorni): 4,8\,\%;
  \item \textbf{moderato} (da 7 a 28 giorni): \textbf{52,9\,\%}, la classe più numerosa;
  \item \textbf{grave} (oltre 28 giorni): \textbf{33,7\,\%};
  \item senza risposta: 8,6\,\%.
\end{itemize}

\rubrica{Situazione e modalità}
\begin{itemize}
  \item \textbf{situazione}: allenamento tecnico-tattico 47,1\,\%, partita 40,4\,\%, allenamento
    fisico 12,5\,\% --- l'allenamento nel suo insieme vale dunque il \textbf{59,6\,\%};
  \item \textbf{modalità}: \textbf{senza contatto} 51,9\,\%, contatto diretto 46,2\,\%, non
    riportato 1,9\,\%;
  \item \textbf{tipo}: primo infortunio 58,6\,\%, \textbf{recidiva} 40,4\,\%.
\end{itemize}

\rubrica{Che cosa se ne ricava}
\begin{itemize}
  \item il \textbf{54\,\%} del campione ha subito almeno un infortunio;
  \item l'\textbf{86,5\,\%} degli infortuni interessa gli \textbf{arti inferiori}, e caviglia più
    ginocchio da soli superano il 50\,\%;
  \item il \textbf{59,6\,\%} avviene \textbf{in allenamento} e il \textbf{51,9\,\%}
    \textbf{senza contatto}: due quote che rimandano entrambe al carico proposto dallo staff, non
    all'avversario;
  \item il peso di caviglia e ginocchio riporta all'importanza degli \textbf{appoggi} in questo
    sport (cap.~\ref{cap:tecnicatattica}), dove il cambio di appoggio è anche mezzo di finta: le
    sollecitazioni articolari che ne derivano vanno sostenute con sedute pensate a quello scopo.
\end{itemize}

% ---------------------------------------------------------------------------
\section{Il quadro d'insieme}

Gli studi visti convergono su una stessa indicazione operativa.

\begin{itemize}
  \item il giocatore di futsal \textbf{si adatta in fretta} a ciò che gli si somministra --- e
    \textbf{solo} a quello: le qualità non sollecitate restano ferme (la velocità in preparazione)
    o \textbf{regrediscono} (il decremento dell'RSA in stagione);
  \item la quota di infortuni \textbf{da sovraccarico} e la quota che matura \textbf{in
    allenamento} indicano che il rischio si genera dove il carico è deciso dallo staff, non dove
    lo impone l'avversario --- e i due studi epidemiologici, su popolazioni del tutto diverse
    (nazionale spagnola maschile, campionato brasiliano femminile), concordano: 47\,\% e
    59,6\,\% degli infortuni in allenamento;
  \item nel \textbf{femminile} il livello di qualificazione non si spiega con i parametri
    fisico-atletici puri, ma con l'\textbf{agility} e la \textbf{tecnica}: un risultato che sposta
    il peso della valutazione verso le prove specifiche;
  \item ne segue che la \textbf{prima azione preventiva} del preparatore, come dell'allenatore, è
    \textbf{far restare il soggetto entro la propria capacità di carico}
    (cap.~\ref{cap:capacitadicarico}) --- e per restarci occorre \textbf{controllare} il carico
    somministrato e \textbf{misurarne gli effetti}.
\end{itemize}

Il controllo dell'allenamento e della partita --- sul piano fisiologico e, se se ne hanno i mezzi,
cinematico --- più i test di valutazione in momenti determinati della stagione servono dunque
\textbf{non solo} al miglioramento della prestazione, ma \textbf{soprattutto} a una vera azione di
\textbf{prevenzione degli infortuni}. In questo senso la capacità di carico
(cap.~\ref{cap:capacitadicarico}) trova qui la sua \textbf{verifica sperimentale}: CK e rapporto
T:C sono gli indicatori con cui si accerta che il carico è stato \textbf{rielaborato} e non subìto.
